\section{Limitations in full}
\label{sec:limitations-full}
\textbf{Substrate.} Acts 2 and 3 measure function-level Python tasks from two code benchmarks, and
Act 4 one scientific-code benchmark, small results artefacts built from it, and four tabular
datasets with synthetic faults. Nothing here establishes that the same ceiling, or the same
explanations, hold for another language, for repository-scale changes, for code whose correctness is
not decided by an executable suite, or for scientific artefacts larger than a model can read whole.
The results and data studies' faults are injected; the code study's failures are the generator's
own, judged by a benchmark that contains defective specifications and tests
\citep{hu2026scicodeverified}. One generator wrote every program. The code and data studies read two
auditor routes, the results study one, each at its default sampling.

\textbf{The ground truth is a test suite, and a test suite is not a specification.} Hidden
suites encode a reference implementation's behaviour, including behaviour its prose never fixed.
In the intervention study reported in \S\ref{sec:inreview}, a generator shown an instance's
failing inputs inferred a behavioural rule that \emph{contradicted} the hidden oracle on 23 of
32 instances. The study cannot say whether that reflects the specifications or the generator's
inference from inputs alone; either way it bounds what any recall figure measured this way can
mean.

\textbf{Interval coverage, measured only in part.} The primary intervals are percentile
bootstraps over problem clusters (over items within dataset in the data study). For a \emph{single rate} only, simulated under the design's
own cluster profiles and an assumed within-problem correlation, with 600 resamples rather than
the 10{,}000 the analyses use (400 replicates; Monte Carlo standard error 0.01--0.02), a nominal
$95\%$ interval covers between $0.93$ and $0.97$ at the tested true rates from $0.10$ to $0.60$, and $0.80$ at a rate of $0.03$ on the
56-cluster stratum (\S\ref{sec:stats}), which suggests that small-rate intervals may be narrower
than they claim. Two earlier interval constructions in this programme were found by review to cover $0.416$ and $0.075$ where $0.95$
was claimed.

\textbf{The validation instrument shares the failure mode under study.} Reports were checked by a
language model, which is what this paper measures a ceiling on. Studies cleared review at between
three and twenty-one rounds, with the round before the last always returning substantive
findings. Measurements moved on review in three studies; elsewhere the refusals concerned
interpretation, where the instrument is weakest. A further pass under an \emph{attempted}
conclusion-suppressed prompt produced six substantive findings that six briefed passes had missed,
but the supplied artefacts exposed the conclusions anyway, so blinding was not achieved and the
comparison is uncontrolled. The defect class it
leaves open is the one that looks right to a language model, and we do not know what remains in
it.

\textbf{The reviewer is one of the measured families.} \texttt{gpt-6-astra} reviewed every
report in this programme and is also an auditor route whose ranking against the others the paper
reports. For claims about \texttt{astra}'s relative performance the validation is a
self-assessment, and those comparisons should be treated as bracketed until a model that audited
nothing here has checked them.

\textbf{No human rater.} Every rating and adjudication in this work was made by a model, one of
them the agent that designed the studies and wrote the rubrics (\S\ref{sec:methods-review}). A
third model rating does not discharge that, and the design that asks for an outside human rater
has not been run.

\textbf{Budget.} Reading depth $K$, the number of arms, and the size of every intervention
population were set by cost rather than by a power calculation, and no sample size was justified
in advance by a minimum-detectable-effect argument.

